% --- Case study box styles (two-color scheme) ---
\definecolor{caseblueframe}{RGB}{67,119,183}
\definecolor{casebluebg}{RGB}{231,239,250}
\definecolor{casegrayframe}{RGB}{120,120,120}
\definecolor{casegraybg}{RGB}{243,243,243}

\newtcolorbox{CasePromptBox}[1]{
  breakable,
  colback=casebluebg,
  colframe=caseblueframe,
  colbacktitle=caseblueframe,
  coltitle=white,
  fonttitle=\bfseries,
  title={#1}
}

\newtcolorbox{CaseTraceBox}[1]{
  breakable,
  colback=casebluebg,
  colframe=caseblueframe,
  colbacktitle=caseblueframe,
  coltitle=white,
  fonttitle=\bfseries,
  title={#1}
}

\newtcolorbox{CaseAnswerBox}[1]{
  breakable,
  colback=casegraybg,
  colframe=casegrayframe,
  colbacktitle=casegrayframe,
  coltitle=white,
  fonttitle=\bfseries,
  title={#1}
}

\setcounter{subsection}{0}
\setcounter{subsubsection}{0}
\renewcommand{\thesubsection}{\arabic{subsection}}
\renewcommand{\thesubsubsection}{\thesubsection.\arabic{subsubsection}}

\subsection{Atomic Tool Alignment}
\label{app:atomic_tool_alignment}

This subsection provides a detailed capability-level alignment for the core atomic tools discussed in the main text.
The purpose is not to enumerate the full runtime tool inventories of Claude Code or OpenClaw, but to show that the core capabilities retained by GA each have corresponding prototypes in both systems.

\begin{table*}[htbp]
\centering
\caption{\textbf{Capability-level alignment of the basic tooling environment.} Each row maps one capability category to the corresponding tool in Claude Code, OpenClaw, and GA.}
\def\arraystretch{0.99}
\setlength{\tabcolsep}{0.42em}
\resizebox{1.0\linewidth}{!}{
\begin{tabular}{llll}
\toprule
\textbf{Capability Category} & \multicolumn{3}{c}{\textbf{Corresponding Tool / Mechanism}} \\
\cmidrule(lr){2-4}
 & \textbf{Claude Code} & \textbf{OpenClaw} & \textbf{GA} \\
\midrule
Code execution & BashTool & exec / process / nodes & code\_run \\
File reading & FileReadTool & read & file\_read \\
File writing & FileWriteTool & write & file\_write \\
Local file editing & FileEditTool & write / patch-like editing & file\_patch \\
Web reading  & WebFetchTool / browser read & browser base reading & web\_scan \\
Web interaction & WebBrowserTool / browser actions & browser interaction & web\_execute\_js \\
Working memory & TodoWriteTool / BriefTool & session history / status-like ability & update\_working\_checkpoint \\
\bottomrule
\end{tabular}
}
\label{tab:appendix_tool_alignment}
\end{table*}

Table~\ref{tab:appendix_tool_replacement_examples} complements the capability-level mapping above with representative task-level replacement examples.
Rather than claiming that GA reproduces every specialized tool one by one, the point is that many benchmarked capabilities can be reconstructed through short compositions of a few atomic tools.

\begin{table*}[t]
\centering
\small
\def\arraystretch{1.02}
\setlength{\tabcolsep}{0.35em}
\begin{tabularx}{\textwidth}{@{}>{\raggedright\arraybackslash}p{0.12\textwidth}>{\raggedright\arraybackslash}p{0.18\textwidth}>{\raggedright\arraybackslash}X>{\raggedright\arraybackslash}X>{\raggedright\arraybackslash}X@{}}
\toprule
\textbf{Source} & \textbf{Specialized Tool / Capability} & \textbf{Representative Task} & \textbf{GA Replacement Composition} & \textbf{Argument} \\
\midrule
Claude Code & \texttt{GlobTool} & \texttt{task\_cc\_01\_glob\_markdown\_list} & \texttt{code\_run} & File-pattern matching can be completed through general code execution. \\
Claude Code & \texttt{GrepTool} & \texttt{task\_cc\_02\_grep\_todo\_hits} & \texttt{code\_run} + \texttt{file\_read} & Text search and localization can be implemented by combining general execution with file reading. \\
\cmidrule(lr){1-5}
Claude Code & \texttt{WebFetchTool} & \texttt{task\_cc\_03\_webfetch\_example\_brief} & \texttt{web\_scan} + \texttt{web\_execute\_js} + \texttt{code\_run} & Web fetching, cleaning, and summarization can be reconstructed by combining browser primitives with code-based post-processing. \\
Claude Code & \texttt{AgentTool} & \texttt{task\_cc\_04\_agent\_fact\_extract} & \texttt{file\_read} + \texttt{code\_run} + \texttt{update\_working\_checkpoint} & Explicit agent delegation can be approximated through stored subagent strategies in memory together with general execution triggers. \\
\cmidrule(lr){1-5}
Claude Code & \texttt{WebSearchTool} & \texttt{task\_cc\_05\_websearch\_example\_results} & \texttt{web\_execute\_js} + \texttt{web\_scan} & Online search can be completed through browser interaction plus webpage reading. \\
Claude Code & \texttt{NotebookEditTool} & \texttt{task\_cc\_06\_notebook\_inspect} & \texttt{file\_read} + \texttt{code\_run} & Notebook-structure inspection can be handled by reading the \texttt{.ipynb} JSON and executing a parsing script. \\
\cmidrule(lr){1-5}
OpenClaw & \texttt{nodes} / structured data & \texttt{task\_oc\_02\_csv\_stats} & \texttt{code\_run} & CSV statistics can be completed through general code execution without a dedicated structured-data tool. \\
\bottomrule
\end{tabularx}
\caption{\textbf{Representative task-level substitutions from specialized tools to GA atomic-tool compositions.} Each row shows one concrete capability test, the original specialized tool or mechanism used by the baseline system, and the corresponding GA composition that can complete the same class of task.}
\label{tab:appendix_tool_replacement_examples}
\end{table*}

\FloatBarrier

\subsection{Web Browsing Visualization}
\label{app:web_browsing_visualization}

Figure~\ref{fig:web_browsing_dual_axis} provides a visual comparison of token consumption and normalized performance across the three web browsing benchmarks discussed in Section~\ref{sec:exp_browsing}.

\begin{figure}[htbp]
\centering
\includegraphics[width=0.7\linewidth]{image/web_browsing_dual_axis.pdf}
\caption{\textbf{Token consumption and normalized score across three web browsing benchmarks.} The left axis shows total token consumption in millions (M), while the right axis shows normalized scores on a 0--1 scale. GA consistently achieves competitive or superior performance while consuming significantly fewer tokens than OpenClaw across all three benchmarks.}
\label{fig:web_browsing_dual_axis}
\end{figure}

\FloatBarrier

\subsection{Case Studies}
\label{app:case_studies}

We present four representative cases corresponding to the main evaluation dimensions.
For readability, the cases are ordered as tool use, memory, self-evolution, and web browsing.
Each case follows a ``purpose and task setup $\rightarrow$ artifacts / trace $\rightarrow$ observed outcome'' pattern so that the appendix reads as a set of worked examples rather than a list of compressed summaries.

\begin{itemize}[leftmargin=*]
\item \textbf{Case A1 (Tool Use).} An API procurement workflow requiring web lookup, structured extraction, budget reasoning, and file generation, used to show how GA handles a long-horizon task with a minimal atomic-tool set (Section~\ref{sec:dimension4_tool_use_efficiency}).

\item \textbf{Case A2 (Memory).} Dangerous-goods hazard classification under condensed, full, and redundant memory variants, used to show why information density matters more than information volume (Section~\ref{sec:exp_memory}).

\item \textbf{Case A3 (Self-Evolution).} A GitHub PR research task showing the progression from natural-language SOP to compiled Python code and a transferability note, used to demonstrate how GA converts experience into reusable executable assets (Section~\ref{sec:self_evolution_capability}).

\item \textbf{Case A4 (Web Browsing).} A BrowseComp-ZH multi-hop reasoning chain where GA identifies a historical figure through iterative search refinement, used to illustrate structured browser extraction in practice (Section~\ref{sec:exp_browsing}).
\end{itemize}

\subsubsection{Case A1. Tool Use: API Procurement Workflow}
\label{app:case_tools}

\begin{CasePromptBox}{Purpose and Task Setup}
\small\ttfamily
Purpose: to examine whether GA's atomic-tool design can support a long-horizon, multi-step workflow, \\
we provide a laboratory procurement task over public LLM pricing pages. \\
This task is useful because it requires browsing, structured extraction, numerical reasoning, and file generation in one workflow. \\
The goal of the case is to show that GA's minimal atomic-tool set can complete the full chain and reach the correct recommendation, providing a qualitative counterpart to Table~\ref{tab:appendix_tool_case}. \\
Task: investigate laboratory procurement for text-only LLM APIs using these public pricing pages: \\
- OpenAI API pricing: https://openai.com/api/pricing/ \\
- Anthropic pricing: https://www.anthropic.com/pricing \\
- Gemini API pricing: https://ai.google.dev/gemini-api/docs/pricing \\
Constraints: \\
- only text models are considered \\
- laboratory size: 6 students \\
- monthly input tokens: 80M \\
- monthly output tokens: 20M \\
- budget: \$300 / month \\
Requirements: \\
1. record candidate flagship text models and their input/output prices \\
2. test whether a single-model plan fits the budget \\
3. if not, trigger fallback and test a dual-model plan \\
4. if that still fails, output infeasible \\
5. final answer must state the primary option, whether it is feasible, whether fallback was triggered, the recommended plan, the estimated monthly cost, and the reason \\
Required output files: cost\_comparison.csv and decision\_04.json.
\end{CasePromptBox}

\begin{CaseTraceBox}{Artifact 1: Cost Comparison File}
\small\ttfamily
cost\_comparison.csv \\
provider,model,input\_cost\_per\_1m,output\_cost\_per\_1m,estimated\_monthly\_cost\_usd,meets\_budget \\
OpenAI, GPT-5.4, 2.5, 15.0, 500.0, No \\
OpenAI, GPT-5.4 mini, 0.75, 4.5, 150.0, Yes \\
OpenAI, GPT-5.4 nano, 0.2, 1.25, 41.0, Yes \\
Anthropic, Claude Opus 4.6, 5.0, 25.0, 900.0, No \\
Anthropic, Claude Sonnet 4.6, 3.0, 15.0, 540.0, No \\
Anthropic, Claude Haiku 4.5, 1.0, 5.0, 180.0, Yes \\
Google, Gemini 2.5 Pro, 1.25, 10.0, 300.0, Yes \\
Google, Gemini 2.5 Flash, 0.3, 2.5, 74.0, Yes \\
Google, Gemini 2.5 Flash-Lite, 0.1, 0.4, 16.0, Yes
\end{CaseTraceBox}

\begin{CaseTraceBox}{Artifact 2: Decision File}
\small\ttfamily
primary\_option: Gemini 2.5 Pro \\
primary\_option\_feasible: true \\
fallback\_triggered: false \\
recommended\_plan: single\_model \\
recommended\_models: [Gemini 2.5 Pro] \\
estimated\_monthly\_cost\_usd: 300.0 \\
reason: Gemini 2.5 Pro is a strong single-model option that still fits within the budget. \\
Cost calculation: 80M input x \$1.25/MTok + 20M output x \$10.00/MTok = \$300.0/month. \\
Other flagship models exceed the budget: GPT-5.4 at \$500/month and Claude Opus 4.6 at \$900/month. \\
No fallback to a dual-model plan is needed.
\end{CaseTraceBox}

\begin{CaseAnswerBox}{Observed Outcome}
\small\ttfamily
GA completed the task in 21 requests and 364,385 total tokens (324.6s). \\
All grading criteria passed: cost\_comparison.csv and decision\_04.json both created with correct schema. \\
Final recommendation: Gemini 2.5 Pro as a single-model plan at exactly \$300/month. \\
Cost breakdown: 80M input $\times$ \$1.25/MTok + 20M output $\times$ \$10.00/MTok = \$300.0. \\
Other flagship models exceeded the budget: GPT-5.4 at \$500/month, Claude Opus 4.6 at \$900/month. \\
No fallback to a dual-model plan was needed.
\end{CaseAnswerBox}

\begin{table}[t]
\centering
\caption{Cross-system comparison on the long-horizon procurement task.}
\label{tab:appendix_tool_case}
\begin{tabular}{lcccc}
\toprule
\textbf{System} & \textbf{Requests} & \textbf{Total Tokens} & \textbf{Time (s)} & \textbf{Success} \\
\midrule
GenericAgent & 21 & 364,385 & 324.64 & Yes \\
Claude Code & 29 & 485,335 & 224.53 & Yes \\
OpenClaw & 12 & 428,402 & 108.42 & Yes \\
\bottomrule
\end{tabular}
\end{table}

\FloatBarrier

\subsubsection{Case A2. Memory: Dangerous-Goods Hazard Classification}
\label{app:case_memory}

\begin{CasePromptBox}{Purpose and Task Setup}
\small\ttfamily
Purpose: to examine whether memory quality depends more on information density than on raw information volume, \\
we provide a dangerous-goods classification task under three memory variants. \\
This task is useful because the decision rule stays fixed while only the memory representation changes. \\
The goal of the case is to show that the main effect comes from how the rules are packaged and retrieved, rather than from whether the rules exist at all. \\
Task: determine hazard\_class from product\_id and the four source texts. \\
Inputs: \\
- SDS label text \\
- handling/storage guidance \\
- transportation requirements \\
- disposal guidance \\
The output must contain both hazard\_score and hazard\_class.
\end{CasePromptBox}

\begin{CaseTraceBox}{Artifact 1: Condensed Memory}
\small\ttfamily
Required rules: \\
1. First validate product\_id. It must be resolvable and follow format P\_XXXXX. \\
\quad If invalid, stop and output hazard\_score=0 and hazard\_class="Unable to Decide". \\
2. Get four component scores: \\
\quad - safety score from SDS label text \\
\quad - handling score from handling/storage guidelines \\
\quad - transportation score from transportation requirements \\
\quad - disposal score from disposal guidelines \\
\quad Each valid component score must be in [1,5]. \\
3. Compute hazard\_score: \\
\quad - If one component is missing or 0, replace it with the max of the other available scores. \\
\quad - If more than two components are missing or 0, output hazard\_score=0 and hazard\_class="Unable to Decide". \\
\quad - Otherwise hazard\_score = sum of the four component scores. \\
\quad - Final valid range is 4 to 20. \\
4. Map hazard\_score to class: \\
\quad - 4-7 to Hazard Class A \\
\quad - 8-12 to Hazard Class B \\
\quad - 13-16 to Hazard Class C \\
\quad - 17-20 to Hazard Class D \\
5. Final output must include both hazard\_score and hazard\_class.
\end{CaseTraceBox}

\begin{CaseTraceBox}{Artifact 2: Full Memory}
\small\ttfamily
1. Purpose \\
To establish a standardized methodology for the systematic identification and classification of dangerous goods hazard classes through multi-source data integration and quantitative severity assessment protocols. \\
2. Scope \\
This procedure encompasses all dangerous goods shipment classification processes within the organization's supply chain operations. \\
3. Definitions \\
SDS = Safety Data Sheet \\
HS = Handling and Storage Guidelines \\
TR = Transportation Requirements \\
DG = Disposal Guidelines \\
AIP = API Integration Protocol \\
HCM = Hazard Classification Matrix \\
SAS = Severity Assessment Score \\
4. Input \\
Product ID (format: P\_XXXXX) \\
Source documentation: SDS / Handling and Storage / Transportation / Disposal \\
API access credentials: endpoint URLs / authentication tokens / backup authentication protocols \\
5. Main Procedure \\
Validate product identification documentation completeness. \\
If product ID fails format requirements, no further action should be taken. \\
The hazard score should be marked as 0 and hazard class as Unable to Decide. \\
Extract four scores, each between 1 and 5. \\
hazard\_score = safety score + handling score + transportation score + disposal score \\
If any component is missing or 0, impute it by taking the max of the other scores. \\
If more than two component scores are missing, mark Unable to Decide. \\
6. Output \\
Final hazard class designation \\
Digital record in Hazard Classification Registry \\
API response logs for all scoring calculations \\
Classification audit trail documentation \\
Final output should be in XML format with tags \textless{}hazard\_score\textgreater{} and \textless{}hazard\_class\textgreater{}
\end{CaseTraceBox}

\begin{CaseTraceBox}{Artifact 3: Redundant Memory}
\small\ttfamily
You are following a standardized dangerous-goods classification procedure designed for supply-chain compliance and multi-source hazard assessment. \\
Definitions: \\
- SDS = safety data sheet \\
- HS = handling and storage guidance \\
- TR = transportation requirements \\
- DG = disposal guidelines \\
- SAS = severity assessment score from 1 to 5 \\
- HCM = hazard classification matrix \\
Input scope: product\_id, SDS text, handling/storage text, transportation text, disposal text, API access and validation context. \\
Decision workflow: \\
1. Validate documentation completeness and product identifier validity. \\
2. product\_id must match format P\_XXXXX and be resolvable. \\
3. If invalid, do not continue. Set hazard\_score to 0 and hazard\_class to ``Unable to Decide''. \\
4. Analyze SDS, handling/storage, transportation, and disposal sources separately. \\
5. Derive one component score per source, each between 1 and 5. \\
6. If exactly one component is missing or 0, impute it using the maximum of the other component scores. \\
7. If more than two components are missing, return ``Unable to Decide'' and hazard\_score 0. \\
8. Sum the four components into a cumulative hazard\_score and validate that total score lies in the range 4-20. \\
9. Convert the total score into Hazard Class A / B / C / D. \\
10. Return the final result in structured form with hazard\_score and hazard\_class, preserving an audit-ready trail. \\
Background note: the broader SOP also mentions registry logging, API integration, source-document handling, and record retention, but the operationally decisive rules are the identifier check, missing-value handling, score summation, and class mapping.
\end{CaseTraceBox}

\begin{CaseAnswerBox}{Observed Outcome}
The memory ablation becomes visually intuitive once the three variants are shown side by side.
The condensed file places every behavior-changing rule in a single short block, whereas the full file spreads the same operational core across purpose statements, definitions, credential notes, registry requirements, and output-format details.
Table~\ref{tab:memory_ablation} can therefore be read primarily as an information-placement result rather than a simple information-volume result.
\end{CaseAnswerBox}

\FloatBarrier

\subsubsection{Case A3. Self-Evolution: GitHub PR Research}
\label{app:case_self_evolution}

\begin{CasePromptBox}{Purpose and Task Setup}
\small\ttfamily
Purpose: to examine whether repeated execution can be distilled into reusable assets, \\
we provide a GitHub PR research task that evolves from an SOP into executable code. \\
This task is useful because the representation shift is directly observable: the agent first stabilizes a procedure as text and then compiles that procedure into code. \\
The goal of the case is to show that what changes across iterations is not merely token count, but the form of reusable knowledge, which serves as the qualitative counterpart of Table~\ref{tab:self_evolution_rounds}. \\
Investigate the five most recent merged bug-fix PRs in a GitHub repository. \\
For each PR: recover the linked issue, identify the affected module, and verify whether the module has troubleshooting coverage in the documentation site. \\
Produce a structured JSON report.
\end{CasePromptBox}

\begin{CaseTraceBox}{Artifact 1: SOP}
\small\ttfamily
Scenario: batch investigation of GitHub PRs, including PR details, linked issues, affected modules, and documentation coverage. \\
Core strategy: \\
DO NOT browse PR pages one by one -- efficiency is extremely low and the run may exceed the turn budget. \\
PREFER batch extraction -- use a script or API to obtain multiple PR records at once. \\
Recommended solution: use github\_pr\_analyzer.py to complete the whole chain in one pass. \\
Example command: \\
python3 ../memory/github\_pr\_analyzer.py langchain-ai/langchain \\
\quad --doc-url https://python.langchain.com \\
\quad --limit 5 \\
\quad --output report.json \\
Script advantages: \\
- pure Python, no browser environment required \\
- one-click execution of the full workflow \\
- flexible command-line parameters \\
- built-in error handling \\
Manual fallback flow: \\
1. build a filtered PR URL \\
2. use document.querySelectorAll(.js-issue-row) to extract PR numbers and titles \\
3. batch fetch PR HTML \\
4. recover issue links with the rule /issues/\textbackslash d+, then Fixes/Closes \#\textbackslash d+, then standalone \#\textbackslash d+ \\
5. infer module names from PR title prefixes such as community: \\
6. verify documentation coverage against troubleshooting and integration paths \\
7. dump the final report with json.dump() \\
Pitfalls: \\
1. do not use window.location.href to visit PR pages one by one \\
2. prefer browser fetch() when manual HTML collection is needed \\
3. distinguish issue links from pull-request links \\
4. different projects may use different documentation subdomains
\end{CaseTraceBox}

\begin{CaseTraceBox}{Artifact 2: Compiled Code}
\small\ttfamily
\textbf{fetch\_pr\_list(...)} \\
\quad url = https://github.com/\{repo\}/pulls?q=\{filters\} \\
\quad response = session.get(url, timeout=30) \\
\quad soup = BeautifulSoup(response.text, 'html.parser') \\
\quad pr\_elements = soup.select('.js-issue-row') \\
\quad collect number, title, and url for the first limit PRs \\
\textbf{extract\_pr\_details(pr)} \\
\quad \detokenize{module_match = re.match(r'^([^:\[]+)', pr['title'])} \\
\quad pattern 1: full GitHub issue URL \\
\quad pattern 2: Fixes|Closes|Resolves \#N \\
\quad pattern 3: standalone \#N excluding the PR number itself \\
\quad return pr\_number, bug\_module, issue\_link \\
\textbf{check\_doc\_coverage(module)} \\
\quad probe multiple candidate paths: \\
\quad /docs/troubleshooting/ \\
\quad /docs/troubleshooting/\{module\} \\
\quad /docs/integrations/\{module\}/ \\
\quad /docs/modules/\{module\}/ \\
\quad /api\_reference/\{module\}/ \\
\quad /docs/integrations/providers/\{module\}/ \\
\quad accept the module if the page contains both the module name and troubleshooting \\
\textbf{analyze(...)} \\
\quad fetch PR list, then extract details, then check documentation, then call json.dump(results, f, indent=2, ensure\_ascii=False)
\end{CaseTraceBox}

\begin{CaseTraceBox}{Artifact 3: Transferability Note}
\small\ttfamily
Original task core abilities: \\
- batch retrieval of GitHub PR lists with filters \\
- extraction of PR details, titles, modules, and linked issues \\
- verification of external documentation coverage \\
- generation of structured JSON reports \\
Reusable asset set: \\
- SOP document: github\_pr\_research\_sop.md \\
- Python script: github\_pr\_analyzer.py \\
Transfer targets described in the case note: \\
1. GitHub issue batch research \\
2. contributor analysis \\
3. release tracking \\
4. workflow and CI inspection \\
5. code search and pattern extraction \\
Example reuse in the issue-analysis case: \\
- keep the batch-fetch strategy \\
- replace /repos/\{repo\}/pulls with /repos/\{repo\}/issues \\
- reuse regex extraction and concurrent request handling \\
Estimated adaptation cost in the transferability note: about 30\% code changes, mainly endpoint replacement and field mapping.
\end{CaseTraceBox}

\begin{CaseAnswerBox}{Observed Outcome}
\small\ttfamily
LangChain verification case (2026-04-08): \\
python3 github\_pr\_analyzer.py langchain-ai/langchain --doc-url https://python.langchain.com --limit 5 \\
\\
5 merged bug-fix PRs retrieved \\
modules found: community (2), partners, openai, core/anthropic \\
issue links recovered: 4/5 \\
troubleshooting coverage found: 4/5 \\
runtime: about 30--40 seconds
\end{CaseAnswerBox}

\FloatBarrier

\subsubsection{Case A4. Web Browsing: BrowseComp-ZH Reasoning Chain}
\label{app:case_browsing}

\begin{CasePromptBox}{Purpose and Task Setup}
\small
Purpose: to examine multi-hop web browsing under a compressed browser interface, \\
we provide a BrowseComp-ZH question whose answer must be recovered through iterative query refinement. \\
This example is useful because the answer is not available from a single lookup and the chain of evidence must be assembled step by step. \\
The goal of the case is to make the narrowing process visible: identify the \emph{shihua} pioneer, confirm Ouyang Xiu, and then recover Wang Anshi from the quotation. \\
Original question (BrowseComp-ZH): \\
``There is a writer A who participated in a literary reform movement, \\
advocated a plain but disciplined style, opposed overly strange and risky writing, \\
and pioneered the \emph{shihua} form of poetic criticism. \\
Another contemporary writer B commented that he had ``deep character'' and ``broad knowledge''. \\
Who was writer B?'' \\
Recorded case metadata: dataset = browsecomp, score = 1.0, time = 130.4s, tokens = input 30050 / output 1296, turns = 7, tools = 6.
\end{CasePromptBox}

\begin{CaseTraceBox}{Artifact 1: Interaction Trace}
\small\ttfamily
Turn 1 summary: start by searching ``first shihua writer'' to identify writer A. \\
Action query: shihua inventor plain style \\
Turn 2 summary: scan the Google result page and narrow the candidate set. \\
Turn 3 summary: because the first search does not uniquely resolve the target, search ``Ouyang Xiu + Liuyi Shihua''. \\
Action query: Ouyang Xiu Liuyi Shihua \\
Turn 4 summary: confirm writer A as Ouyang Xiu. \\
Turn 5 summary: search the evaluative phrase together with Ouyang Xiu. \\
Action query: deep character broad knowledge Ouyang Xiu \\
Turn 6 summary: verify that the quoted evaluator is Wang Anshi. \\
Turn 7 summary: return the final answer with confidence.
\end{CaseTraceBox}

\begin{CaseTraceBox}{Artifact 2: Final Reasoning Excerpt}
\small\ttfamily
Explanation: \\
Through the searches on ``first shihua writer'' and related keywords, the agent identified writer A as Ouyang Xiu. \\
He participated in the Northern Song literary reform movement, advocated a plain but disciplined style, \\
opposed strange and risky writing, and pioneered the form through Liuyi Shihua. \\
After searching the evaluative phrase together with Ouyang Xiu, \\
the agent found that Wang Anshi, in a commemorative text for Ouyang Xiu, \\
described him as having ``deep character'' and ``broad knowledge''. \\
Therefore the writer B in the question is Wang Anshi.
\end{CaseTraceBox}

\begin{CaseAnswerBox}{Observed Outcome}
\small\ttfamily
score: 1.0 \\
turns: 7 \\
tool calls: 6 \\
time: 130.4s \\
Exact Answer: Wang Anshi \\
Judge result: correct = yes, confidence = 95\% \\
Reasoning note: the extracted final answer exactly matches the gold answer ``Wang Anshi''.
\end{CaseAnswerBox}

\FloatBarrier

\subsection{General Capability Showcase}
\label{app:general_capability_showcase}

Beyond the four evaluation-aligned cases above, we present five additional cases drawn from a pool of more than 500 historical sessions.
They are selected to maximize coverage of non-benchmark capability dimensions rather than to reflect average-case performance, and are intended as qualitative demonstrations rather than statistical evaluation.
All cases come from user-authorized sessions; sensitive identities, document contents, account details, and message contents are anonymized or paraphrased when presented here, and any autonomous behaviors were executed under user-configured triggers or preset operating policies.
Together, these cases illustrate how GA's architectural choices --- atomic tools, persistent memory, and self-evolution --- support real-world workflows that are difficult to capture in standardized benchmarks.
Each case follows the same ``purpose and task setup $\rightarrow$ execution trace $\rightarrow$ observed outcome'' format for consistency.

\begin{itemize}[leftmargin=*]
\item \textbf{Case B1 (Cross-Device Control).} A mobile food-ordering task via ADB, followed by screen-recording retrieval and video post-processing, used to show how atomic-tool composition can bridge the PC--smartphone boundary within one execution chain.

\item \textbf{Case B2 (Cross-Platform Orchestration).} A message-forwarding task from a local WeChat database to a Weibo post, used to show how GA composes heterogeneous local and web environments within a single workflow.

\item \textbf{Case B3 (Autonomous Operation).} An overnight session triggered after the user leaves, used to show how self-evolution, working checkpoints, and preset patrol routines can support bounded autonomous operation.

\item \textbf{Case B4 (Remote Infrastructure).} An SSH-based remote file-server deployment task, used to show how atomic tools can support end-to-end DevOps workflows including dependency installation, file transfer, service deployment, and iterative troubleshooting.

\item \textbf{Case B5 (Long-Horizon Academic Workflow).} A multi-session NSFC grant proposal assistance task spanning figure design, citation verification, and batch error correction, used to show how persistent memory supports extended, multi-phase academic work.
\end{itemize}

\subsubsection{Case B1. Cross-Device Control: Mobile Food Ordering via ADB}
\label{app:case_cross_device}

\begin{CasePromptBox}{Purpose and Task Setup}
\small\ttfamily
Purpose: to examine whether GA can operate across the PC--smartphone boundary through ADB-based device control, \\
we provide a real-world mobile food-ordering task followed by screen-recording retrieval and video post-processing. \\
This task is useful because it requires coordinated control of a physical mobile device, GUI interaction with a commercial app, and multimedia processing on the host PC within a single workflow. \\
The goal of the case is to show that a small atomic-tool set can be composed into a cross-device execution chain under a realistic consumer-app workflow. \\[6pt]
User instruction (verbatim, translated): \\
``Order two cups of milk tea on my phone via Meituan Waimai. Stop at the final payment page --- do not pay.'' \\[4pt]
Follow-up instructions: \\
1. ``Pull the screen recording from the phone.'' \\
2. ``Trim everything before 1:10, black out the top half after 3:40.'' \\
3. ``Convert it into a sped-up GIF.''
\end{CasePromptBox}

\begin{CaseTraceBox}{Artifact 1: Mobile Ordering Execution Trace}
\small\ttfamily
Step 1: Establish ADB connection to the Android device. \\
Step 2: Launch the Meituan Waimai (food delivery) app via ADB shell am start. \\
Step 3: Dismiss pop-up advertisements by locating and tapping the close button coordinates. \\
Step 4: Navigate to the ``Desserts \& Drinks'' category. \\
Step 5: Select a nearby store (Hushang Ayi, a popular milk-tea chain). \\
Step 6: Locate the ``Must-Try Milk Tea --- Pick Any Two'' combo at \textyen 23.9. \\
Step 7: Select flavor 1: Thick Taro Paste; select flavor 2: Brown Sugar Boba. \\
Step 8: Confirm selections and add to cart. \\
Step 9: Proceed to the checkout page and halt --- no payment action taken.
\end{CaseTraceBox}

\begin{CaseTraceBox}{Artifact 2: Video Post-Processing Trace}
\small\ttfamily
Step 1: Pull the screen recording from the phone via adb pull. \\
Step 2: Trim the first 1 minute 10 seconds using ffmpeg -ss 00:01:10. \\
Step 3: Black out the top half of the frame after the 3:40 mark using ffmpeg drawbox filter with overlay logic. \\
Step 4: Re-encode and speed up the video (4x) to produce a compact GIF using ffmpeg with palettegen and paletteuse filters. \\
Step 5: Verify the final GIF file size and frame count.
\end{CaseTraceBox}

\begin{CaseAnswerBox}{Observed Outcome}
\small\ttfamily
In this session, GA completed the full pipeline: ADB device connection, commercial app navigation (7 GUI interactions), screen-recording retrieval, video trimming, region-specific masking, and GIF generation. \\
The workflow was executed through a composition of code\_run (ADB commands, ffmpeg), file\_read / file\_write (media files), and update\_working\_checkpoint (tracking multi-phase progress). \\
No task-specific mobile-automation framework (e.g., Appium, UIAutomator) was introduced in this workflow. \\
This case specifically illustrates the paper's claim that a compact atomic-tool set can be composed into a long execution chain that extends beyond the host PC when ADB is available as a bridge.
\end{CaseAnswerBox}

\FloatBarrier

\subsubsection{Case B2. Cross-Platform Orchestration: WeChat to Weibo Message Forwarding}
\label{app:case_cross_platform}

\begin{CasePromptBox}{Purpose and Task Setup}
\small\ttfamily
Purpose: to examine whether GA can orchestrate data flow across heterogeneous local and web platforms, \\
we provide a task that requires reading from a local encrypted database and publishing to a social-media platform via browser automation. \\
This task is useful because it chains together local database decryption, contact resolution, message extraction, and browser-based content publishing --- capabilities that belong to different software stacks. \\
The goal of the case is to show that GA can compose heterogeneous data sources and output channels into a single workflow under explicit user instruction. \\[6pt]
User instruction (verbatim, translated): \\
``Forward the latest message from Professor Xiao in WeChat to Weibo.'' \\[4pt]
Presentation note: contact identity and forwarded message content are anonymized here, and the operation was performed only after the user explicitly requested the forwarding action.
\end{CasePromptBox}

\begin{CaseTraceBox}{Artifact 1: Execution Trace}
\small\ttfamily
Phase 1 --- WeChat Message Extraction: \\
Step 1: Locate the local WeChat database files (EnMicroMsg.db). \\
Step 2: Decrypt the database using SQLCipher with the derived key. \\
Step 3: Query the contact table to resolve the anonymized alias ``Professor Xiao'' to the intended contact record. \\
Step 4: Query the message table for the most recent private message from this contact. \\
Step 5: Extract the latest message content (paraphrased here for privacy). \\[4pt]
Phase 2 --- Weibo Publishing: \\
Step 6: Open the Weibo compose page in the browser via web\_scan. \\
Step 7: Inject JavaScript to locate the post editor textarea and fill in the extracted message. \\
Step 8: Trigger the submit button via web\_execute\_js. \\
Step 9: Verify that the post appears on the user's Weibo timeline.
\end{CaseTraceBox}

\begin{CaseAnswerBox}{Observed Outcome}
\small\ttfamily
In this session, GA bridged two heterogeneous platforms within one workflow: \\
(1) local encrypted WeChat database $\rightarrow$ SQLCipher decryption $\rightarrow$ contact resolution $\rightarrow$ message extraction; \\
(2) browser-based Weibo $\rightarrow$ JS injection $\rightarrow$ post creation $\rightarrow$ publication verification. \\
The tools used were code\_run (SQLCipher operations, key derivation), file\_read (database file access), and web\_scan + web\_execute\_js (Weibo interaction). \\
This case specifically illustrates the paper's claim that atomic tools can be composed across heterogeneous local and web environments, allowing one session to move from private local data handling to browser-side execution under explicit user intent.
\end{CaseAnswerBox}

\FloatBarrier

\subsubsection{Case B3. Autonomous Operation: Unsupervised Overnight Session}
\label{app:case_autonomous}

\begin{CasePromptBox}{Purpose and Task Setup}
\small\ttfamily
Purpose: to examine whether GA can operate productively with limited oversight under a preset autonomy policy, \\
we present an autonomous session triggered by the system detecting that the user has been away for over 30 minutes. \\
This task is useful because it demonstrates self-directed task selection, execution, and self-correction over an extended period. \\
The goal of the case is to show that GA's self-evolution capability can extend beyond benchmark settings into bounded autonomous operation. \\[6pt]
Trigger condition: \\
{[AUTO]} User has been away for more than 30 minutes. \\
Operational boundary: the session was restricted to inspection, reporting, tool creation, and reversible local maintenance actions; external publication, payment, and destructive system modification were excluded by policy. \\
No explicit task instruction was given after the trigger; GA selected and executed tasks based on its built-in patrol and self-improvement routines within the preset policy.
\end{CasePromptBox}

\begin{CaseTraceBox}{Artifact 1: Autonomous Task Inventory (Rounds 125--140+)}
\small\ttfamily
Category 1 --- System Security Audit: \\
- Scanned all listening ports on the host machine. \\
- Discovered a stock-trading application occupying 939 UDP ports and 509\,MB of memory at 3:00\,AM. \\
- Audited startup items: 11 registry entries + 267 scheduled tasks + 86 auto-start services. \\[4pt]
Category 2 --- Tool Creation: \\
- Wrote port\_monitor.py: real-time port monitoring utility. \\
- Wrote process\_watchdog.py: process resource watchdog with alerting. \\
- Wrote startup\_auditor.py: startup-item auditing and reporting tool. \\[4pt]
Category 3 --- Web Patrol: \\
- Visited tech community sites (Guohe Boke, V2EX) and GitHub Trending. \\
- Generated a summarized patrol report covering trending topics and notable repositories. \\[4pt]
Category 4 --- Environment Hygiene: \\
- Audited the Python environment: found 294 installed packages occupying 1.9\,GB. \\
- Identified 2 CVE vulnerabilities in the filelock package; wrote pip\_audit.py. \\
- Discovered and fixed an inconsistency in its own history filenames (.md vs .txt extension mismatch).
\end{CaseTraceBox}

\begin{CaseAnswerBox}{Observed Outcome}
\small\ttfamily
Over approximately 15+ autonomous rounds, GA performed system security auditing, created three reusable utility scripts, conducted web patrol with report generation, audited the Python environment for vulnerabilities, and self-corrected its own file-management errors. \\
After the trigger fired, the session proceeded without additional user input and remained within the preset autonomy policy described above. \\
Tools used: code\_run (system commands, script creation and testing), file\_read / file\_write (script files, reports), web\_scan (web patrol), and update\_working\_checkpoint (tracking progress across rounds). \\
This case specifically illustrates the paper's claim that self-evolution, working checkpoints, and patrol routines can support bounded autonomous operation over extended periods, rather than only repeated benchmark-style tasks.
\end{CaseAnswerBox}

\FloatBarrier

\subsubsection{Case B4. Remote Infrastructure: SSH-Based File Server Deployment}
\label{app:case_remote_infra}

\begin{CasePromptBox}{Purpose and Task Setup}
\small\ttfamily
Purpose: to examine whether GA can manage end-to-end remote infrastructure tasks, \\
we provide a real-world deployment scenario requiring SSH access, file transfer, and iterative service configuration. \\
This task is useful because it involves dependency management on an unfamiliar remote machine, multi-step troubleshooting, and responding to changing user requirements mid-task. \\
The goal of the case is to show that an atomic-tool workflow can support a complete DevOps-style execution chain in this setting, without relying on infrastructure-as-code tooling. \\[6pt]
User instruction (verbatim, translated): \\
``SSH into user@{\it [IP address]}. Set up a file server on it and upload the zip file.'' \\
Follow-up instruction: ``Make it publicly accessible.''
\end{CasePromptBox}

\begin{CaseTraceBox}{Artifact 1: Execution Trace}
\small\ttfamily
Phase 1 --- Environment Setup: \\
Step 1: Install the paramiko library for SSH connectivity (pip install paramiko). \\
Step 2: Establish an SSH connection to the remote Linux server. \\
Step 3: Upload the 19\,MB zip file via SFTP. \\[4pt]
Phase 2 --- Service Deployment: \\
Step 4: Deploy a Python HTTP file server on the remote machine (python3 -m http.server). \\
Step 5: Encounter and fix a Chinese filename encoding issue (UTF-8 locale configuration). \\[4pt]
Phase 3 --- Requirement Change: \\
Step 6: User requests public access $\rightarrow$ reconfigure the server to remove authentication. \\
Step 7: Add directory isolation so that only the target folder is exposed. \\[4pt]
Phase 4 --- Verification: \\
Step 8: Test the download from the local machine. \\
Step 9: Confirm file integrity: downloaded size = 19.01\,MB, matching the original.
\end{CaseTraceBox}

\begin{CaseAnswerBox}{Observed Outcome}
\small\ttfamily
In this session, GA completed the full DevOps workflow: SSH connection, dependency installation on the remote machine, file upload (19\,MB), HTTP server deployment, encoding bugfix, mid-task requirement change (public access), directory isolation, and end-to-end verification. \\
Tools used: code\_run (paramiko SSH, SFTP, remote shell commands), file\_read / file\_write (local zip file, configuration), and update\_working\_checkpoint (phase tracking). \\
This case specifically illustrates the paper's claim that atomic tools plus lightweight state tracking can support a complete ``connect $\rightarrow$ configure $\rightarrow$ deploy $\rightarrow$ troubleshoot $\rightarrow$ verify'' workflow over SSH in a realistic remote environment.
\end{CaseAnswerBox}

\FloatBarrier

\subsubsection{Case B5. Long-Horizon Academic Workflow: NSFC Grant Proposal Assistance}
\label{app:case_academic}

\begin{CasePromptBox}{Purpose and Task Setup}
\small\ttfamily
Purpose: to examine whether GA can sustain coherent support over a long-horizon, multi-session academic task, \\
we present a real-world NSFC (National Natural Science Foundation of China) grant proposal assistance case that spanned multiple sessions over several days. \\
This task is useful because it involves reading and comprehending a full-length proposal, identifying structural weaknesses, producing academic figures, performing large-scale citation verification, and iteratively correcting errors --- capabilities that must be coordinated across sessions with persistent memory. \\
The goal of the case is to show that GA's memory system and self-evolution capabilities can support sustained, multi-phase academic assistance beyond a single session. \\[6pt]
Presentation note: the proposal materials were user-provided, and document-specific details are anonymized here when they are not essential to the technical point. \\
Task phases (emerged organically from user interaction): \\
1. Read and analyze the complete NSFC proposal document. \\
2. Identify that the proposal contains no figures; design a figure plan. \\
3. Generate academic figures (overview diagram + per-section illustrations). \\
4. Verify all bibliographic entries against external sources. \\
5. Correct discovered citation errors and regenerate the PDF.
\end{CasePromptBox}

\begin{CaseTraceBox}{Artifact 1: Figure Design and Generation}
\small\ttfamily
Observation: the entire proposal contained zero figures, which is a significant weakness for an NSFC application. \\
Action plan: \\
- Design a research overview figure showing the relationship among the three proposed research themes. \\
- Design per-section illustrations for each research content block. \\[4pt]
Execution: \\
- Used Python (matplotlib, networkx) to generate vector-format overview diagrams. \\
- Composed Gemini API prompts to generate domain-appropriate academic illustrations. \\
- Integrated figures into the LaTeX source and verified rendering.
\end{CaseTraceBox}

\begin{CaseTraceBox}{Artifact 2: Citation Verification and Batch Correction}
\small\ttfamily
Approach: parsed the .bib file and extracted all bibliographic entries. \\
For each entry: \\
- Queried arXiv API to verify title, authors, and year. \\
- Cross-referenced with OpenAlex API for published-venue entries. \\
- Compared extracted metadata against the .bib fields. \\[4pt]
Findings: \\
- Discovered a significant number of arXiv preprint entries with incorrect metadata (wrong year, mismatched titles, incomplete author lists). \\
- Root cause: likely auto-generated .bib entries from LLM-assisted writing without manual verification. \\[4pt]
Correction: \\
- Batch-updated the .bib file with verified metadata. \\
- Logged all changes for the user's review (change log with before/after comparisons). \\
- Regenerated the PDF with corrected references.
\end{CaseTraceBox}

\begin{CaseAnswerBox}{Observed Outcome}
\small\ttfamily
Across multiple sessions spanning several days, GA provided end-to-end assistance for: \\
(1) full-document comprehension and structural analysis; \\
(2) figure design and generation (Python vector graphics + API-generated illustrations); \\
(3) systematic citation verification against arXiv and OpenAlex, discovering and correcting a batch of erroneous entries; \\
(4) LaTeX integration and PDF regeneration. \\[4pt]
Tools used: file\_read (proposal document, .bib file), code\_run (Python plotting, API queries, .bib parsing), file\_write (corrected .bib, figures), web\_scan (arXiv / OpenAlex verification), and update\_working\_checkpoint (cross-session state). \\[4pt]
This case specifically illustrates the paper's claim that layered persistent memory (L0--L3), together with reusable procedures accumulated through self-evolution, can maintain continuity across multi-session academic workflows involving reading, generation, verification, and revision.
\end{CaseAnswerBox}

\FloatBarrier
